\section{MDP 把决策问题写成随机过程}
\label{sec:16-Control-and-Reinforcement-Learning/02-Reinforcement-Learning/01}

一台仓储机器人停在通道口。左边那条路近两米，但叉车常堵；右边稳当，要多绕十几米；原地等三秒也许正好让运输车先过。你写调度程序时会发现，比较这三个选项根本不够——选左边不只是省了两米，它还把机器人送到一个新位置，那里能选的动作、会碰到的拥堵，全都变了。评价一个动作，你得连它之后的处境一起评价。

这一节的任务是把这条因果链写成一个可以计算的对象。写法本身不难，难的是里面藏着一个容易被跳过的前提：\textbf{状态必须记住足够多的东西}。这个前提一旦不成立，后面几节的所有算法都会在数学上失去依据，而它们照样会跑完、照样会输出一张漂亮的价值表。

\subsection{一次交互只需要五个对象}

一个折扣 Markov 决策过程（MDP）记作

\[
(\mathcal S,\mathcal A,P,R,\gamma),
\]

其中五个成分各司其职：

\begin{itemize}
\tightlist
\item
  状态空间 \(\mathcal S\)：做当前决策所需的全部信息；
\item
  动作空间 \(\mathcal A\)：可以施加的输入；
\item
  转移核 \(P(s'\given s,a)\)：在 \(s\) 执行 \(a\) 后落到 \(s'\) 的概率；
\item
  奖励 \(R\)：一次转移带来的即时评价，写成 \(R(s,a,s')\) 或它的条件期望 \(r(s,a)\)；
\item
  折扣率 \(0\le\gamma<1\)：未来奖励折回当前的权重。
\end{itemize}

交互在离散时间上循环推进：从 \(S_t\) 出发选 \(A_t\)，环境返回 \(R_{t+1}\) 与 \(S_{t+1}\)，再从 \(S_{t+1}\) 继续。奖励下标写 \(t+1\)，是因为它由时刻 \(t\) 的状态与动作触发，与下一状态一起返回。有的文献写 \(R_t\)。两种约定都能自洽，但你换一本书读之前得先看一眼下标习惯，否则推导里会凭空多出或少掉一个 \(\gamma\)。

拿掉动作，转移只剩 \(P(s'\given s)\)，MDP 就退化为\hyperref[sec:06-Random-Processes/03-Markov-Chains/01]{``Markov 性质、转移矩阵与多步演化''}里讲的那条 Markov 链。动作的加入让转移矩阵不再固定：每选一个动作，你就选走了下一步概率规律的一部分。

\subsection{Markov 性质是对状态设计提的要求}

MDP 要求

\[
P(S_{t+1},R_{t+1}\given S_0,A_0,\ldots,S_t,A_t)
=
P(S_{t+1},R_{t+1}\given S_t,A_t).
\]

这句话常被读成``世界没有记忆''，读错了。世界当然有记忆。这个等式是在说：你选的那个 \(S_t\)，已经把历史中所有还会影响未来的部分\emph{装进去了}。它是对状态变量的一项设计指标，不是对物理世界的一句描述。

恒温控制最容易踩这个坑。假设状态只记当前温度。传感器读到 23 度，你要决定是否继续加热。但``刚关掉加热器''的 23 度和``刚打开加热器''的 23 度，接下来五分钟的走向完全相反——加热器铁芯里的余热还在往房间里放，或者刚开始吸功还没热起来。同一个状态标签下面压着两种不同的未来，Markov 性质在这里已经断了。

\begin{center}
\begin{tikzpicture}[>=stealth, scale=0.92]
% 左：粗状态
\node at (2.4,3.15) {\small 状态只记温度};
\node[draw,rounded corners,minimum width=1.5cm,minimum height=0.75cm] (c) at (2.6,1.3) {\small \(23^\circ\)};
\node (a1) at (0.4,2.3) {\small 刚关加热};
\node (a2) at (0.4,0.3) {\small 刚开加热};
\draw[->] (a1) -- (c);
\draw[->] (a2) -- (c);
\draw[->,dashed] (c) -- (4.4,2.3);
\draw[->,dashed] (c) -- (4.4,0.3);
\node[anchor=west] at (4.5,2.3) {\small 继续升};
\node[anchor=west] at (4.5,0.3) {\small 转降};
% 右：细化状态
\begin{scope}[shift={(7.2,0)}]
\node at (2.4,3.15) {\small 状态加记加热器};
\node[draw,rounded corners,minimum width=1.6cm,minimum height=0.7cm] (d1) at (2.6,2.3) {\small \(23^\circ\)/关};
\node[draw,rounded corners,minimum width=1.6cm,minimum height=0.7cm] (d2) at (2.6,0.3) {\small \(23^\circ\)/开};
\node (b1) at (0.4,2.3) {\small 刚关加热};
\node (b2) at (0.4,0.3) {\small 刚开加热};
\draw[->] (b1) -- (d1);
\draw[->] (b2) -- (d2);
\draw[->,dashed] (d1) -- (4.4,2.3);
\draw[->,dashed] (d2) -- (4.4,0.3);
\node[anchor=west] at (4.5,2.3) {\small 继续升};
\node[anchor=west] at (4.5,0.3) {\small 转降};
\end{scope}
\end{tikzpicture}
\end{center}

\emph{左边两段历史被同一个状态标签合并，未来却分道扬镳；右边把加热器状态并进去之后，每个状态各自决定一个未来。}

图里左右两边的差别，就是 Markov 性质成立与不成立的差别。左边那个 \(23^\circ\) 节点画出去两条虚线，说明``给定状态''之后未来仍然要看更早的历史；右边每个节点只出一条，历史的信息已经被吸收进状态标签本身。检验的操作方式也就写在图里：固定你选的状态，问更早的观测是否还携带关于下一状态或未来回报的额外信息。如果还有，状态就不够。

修法是把还在起作用的那部分历史塞进来——加热器内部温度、上一个动作、最近几步的滑动摘要。这类修法会让状态空间变大，代价真实存在，但比起在一个不成立的递归上跑一万轮迭代，这个代价便宜得多。

若真实状态压根无法完整观测，只看得到观测 \(O_t\)，问题就落到部分可观测 MDP（POMDP）里。此时可以拿整段观测历史、循环网络的隐状态或者对隐藏状态的信念分布当作决策依据。把非 Markov 的观测直接当状态用，是最省事也最危险的做法：算法不会报错，只会安静地收敛到一个错误问题的解。

\subsection{策略一固定，MDP 就退回一条 Markov 链}

策略 \(\pi(a\given s)\) 给出在状态 \(s\) 选动作 \(a\) 的概率；确定性策略是它的特例，写成 \(a=\pi(s)\)。策略固定之后，动作层面的不确定性被吸收进转移：

\[
P^\pi(s'\given s)
=
\sum_a \pi(a\given s)P(s'\given s,a).
\]

受控过程重新变回一条普通 Markov 链。这个化归很实用——评价一个策略，等价于研究它诱导出的那条链的长程行为，于是\hyperref[sec:06-Random-Processes/03-Markov-Chains/04]{``平稳分布、可逆性与遍历收敛''}里的全部工具都可以直接搬过来用。后面几节谈到``访问分布''时，指的就是这条链的分布。

只依赖当前状态、不显式依赖时刻的策略叫平稳 Markov 策略。对有限状态动作的无限时域折扣 MDP，存在确定性的平稳最优策略。这个结论不是自明的，需要证明，但确实成立，也是后面所有算法只在状态上做查表或拟合的理由。有限时域问题就不同了：剩余步数会改变最优动作，你要么把时间写进策略，要么把时间并进状态。把无限时域的结论不加说明地搬到有限时域任务上，是一类常见且不容易察觉的错误。

\subsection{折扣把无限长的后果折回当前一步}

只最大化下一步奖励的算法一定短视。定义从时刻 \(t\) 起的折扣回报

\[
G_t
=R_{t+1}+\gamma R_{t+2}+\gamma^2R_{t+3}+\cdots
=\sum_{k=0}^{\infty}\gamma^kR_{t+k+1}.
\]

睡前纠结``再看一集''还是``现在睡''，就是这个结构：多看一集的放松立刻兑现，少睡的疲劳要到明早才结账。评价标准如果只停在下一步，明早的困倦根本进不了今晚的决策。\(G_t\) 把此后发生的一切沿时间轴折叠回当前，折扣率 \(\gamma\) 决定折叠时每一层压多扁。

奖励有界且 \(\gamma<1\) 时，这个级数绝对收敛，理由与\hyperref[sec:02-Calculus/05-Sequences-and-Series/02]{``无穷级数与部分和''}里的几何级数完全一样。\(\gamma\) 有三种读法，彼此不冲突：它可以是时间偏好，越早的奖励越要紧；可以是随机终止，每步以概率 \(1-\gamma\) 结束，于是 \(\gamma^k\) 正好是任务活过 \(k\) 步的概率；也可以纯粹当作数学正则化，它让下一节的 Bellman 算子成为压缩映射。

要紧的是别把 \(\gamma\) 当成一个只影响速度的技术旋钮。把它调小，最优策略本身会变——机器人会真的更偏好近处的收益，而不是``更快地找到同一个答案''。调到接近 1，有效规划视界拉长，远期奖励能传回来，但远期的估计误差也一起传回来了，收敛还会明显变慢。下一节会把这个取舍量化成一条几何收敛速率。对自然终止且期望长度有限的分幕任务，可以取 \(\gamma=1\)；对不终止的持续任务，常改用长期平均奖励，那时理论和算法都得重新表述一遍。

\subsection{两个值函数回答两种不同的问题}

策略 \(\pi\) 的状态值是

\[
V^\pi(s)
=
\E_\pi\left[G_t\given S_t=s\right],
\]

读作``走到 \(s\) 之后一直按 \(\pi\) 行动，预期能拿多少''。动作值是

\[
Q^\pi(s,a)
=
\E_\pi\left[G_t\given S_t=s,A_t=a\right],
\]

读作``当前强制先做 \(a\)，之后再按 \(\pi\) 行动''。两者由策略加权相连：

\[
V^\pi(s)=\sum_a\pi(a\given s)Q^\pi(s,a).
\]

这两个量都是条件期望，也就继承了\hyperref[sec:05-Probability/04-Expectation-and-Moments/04]{``条件期望是利用信息的最佳预测''}里的全部性质，包括塔性质——下一节的递归展开靠的就是它。

分工上，\(V\) 适合评价处境，\(Q\) 适合直接比较选择。模型已知时，从 \(V\) 做一步前瞻就能比较动作；模型未知时，\(Q\) 的好处是选动作根本不需要显式的转移概率，\(\argmax_a Q(s,a)\) 就够了。第四节的 Q 学习正是从这个缺口进去的。

\subsection{奖励定义了什么叫最优，而不只是加速学习}

仓储机器人若每走一步得 \(-1\)、到达目标得 \(0\)，最大化回报就等价于偏好更短路径。若碰撞罚得不够重，它会去走高风险捷径，而且这不是它学错了——按你写的奖励，那条捷径确实更优。若只在到达时给正奖励、不给步数代价，在 \(\gamma=1\) 下所有最终能到达的路线价值相同，算法在其中随便挑一条也完全合规。

奖励塑形能缓解稀疏反馈，也能悄悄改写任务。按``离目标的距离减少了多少''发奖励，智能体可能在两个位置之间来回刷分；只奖励点击率，推荐系统会牺牲长期满意度去换短期指标。有一类形式是安全的：势函数塑形

\[
F(s,a,s')=\gamma\Phi(s')-\Phi(s)
\]

在标准条件下不改变最优策略集合，因为它沿任何轨迹的折扣累加都退化成一个只依赖起点终点的伸缩项。除此之外，任意加减奖励项都可能重写问题。``学得更快''和``目标没变''是两件事，做实验时要分开报告。

\subsection{模型已知还是只能采样，分界在期望怎么算}

如果 \(P\) 和期望奖励都已知，你可以对所有下一状态精确求和，直接去解下一节的 Bellman 方程，这条路叫动态规划，它与\hyperref[sec:16-Control-and-Reinforcement-Learning/01-Optimal-Control/03]{``HJB 方程与动态规划''}是同一套思想的离散版本。如果它们未知但环境可以交互，你能拿到的只有一串四元组 \((S_t,A_t,R_{t+1},S_{t+1})\)，只能用样本平均或随机逼近去替代那个期望，这才进入强化学习。两条路优化的是同一个 MDP，差别在信息条件，不在目标。

``可以观察单步样本''这句话本身也需要审。设 \(I_t\) 是现有库存，\(q_t\) 是补货量，潜在需求为 \(D_t\)，下一期库存满足

\[
I_{t+1}=(I_t+q_t-D_t)_+ ,
\]

但数据库里记下的只有销量

\[
S_t=\min\set{D_t,\;I_t+q_t}.
\]

补货量 \(q_t\) 不只改变下一状态，它还决定了缺货时有多少需求根本没被记录。如果状态里没有缺货标志、在途订单和促销信息，你采到的转移样本就撑不起你写下的那个 Markov 模型。给算法贴上 model-free 的标签不会让删失需求消失——它只表示更新时没有查一张显式的转移表，不表示你可以不追问数据是怎么来的。

MDP 还默认转移和奖励规律在学习期间大体稳定。若环境因为别的学习者改了策略、因为季节、因为设备磨损而持续漂移，旧样本与当前规律就不再同分布，Bellman 方程本身在移动，多训几轮修不了这件事。稳定性与状态充分性是同一类前提的两面：前者要求规律不随时间变，后者要求状态装得下历史。下一节假定这两条都成立，去解那个已经写好的方程。
